\begin{solution}
The null travel-time functional contains
\((1-(1+\gamma)\Phi)\dd\ell\).  Varying the path gives
\[
\Delta\theta=(1+\gamma)\int_{-\infty}^{\infty}
\frac{Mb\,\dd x}{(b^2+x^2)^{3/2}}
=\frac{2(1+\gamma)M}{b}.
\]
General relativity has \(\gamma=1\), yielding \(4M/b\).
\end{solution}

\begin{solution}
Along the unperturbed ray,
\[
\Delta t=-(1+\gamma)\int\Phi\,\dd\ell.
\]
For endpoints far from closest approach,
\[
\Delta t\simeq(1+\gamma)M
\ln\frac{4r_1r_2}{b^2}.
\]
Thus radar echo delay, like light deflection, primarily tests
\(1+\gamma\).
\end{solution}

\begin{solution}
With vector potential
\(\bm A=\bm J\times\bm r/r^3\), the precession is
\[
\bm\Omega_{\rm LT}
=\frac{1}{r^3}\left[3(\bm J\cdot\hat{\bm r})\hat{\bm r}-\bm J\right].
\]
On the axis it is \(2\bm J/r^3\); in the equatorial plane it is
\(-\bm J/r^3\).  The sign refers to the spin's rotation relative to an
asymptotically inertial frame.
\end{solution}

\begin{solution}
Expanding parallel transport in the weak Schwarzschild metric gives
\[
\bm\Omega_{\rm dS}=\frac32\,\bm v\times\bm\nabla\Phi.
\]
For a circular orbit, integration over \(2\pi r/v\) gives
\(\Delta\phi=3\pi M/r\).  One third of the coefficient is the kinematic
Thomas term and two thirds arise from spatial curvature.
\end{solution}

\begin{solution}
The first post-Newtonian radial equation contains an effective
\(3Mu^2\) correction for \(u=1/r\).  Substituting
\(u_0=[1+e\cos\phi]/[a(1-e^2)]\) and retaining the resonant forcing shifts
the radial frequency, giving
\[
\Delta\varpi=\frac{6\pi M}{a(1-e^2)}.
\]
At this order the relative orbit depends on total mass in the same universal
combination; mass-ratio dependence enters other 1PN observables and higher
orders.
\end{solution}

\begin{solution}
The observed semiamplitude is
\[
K_1=\frac{2\pi a_1\sin i}{P\sqrt{1-e^2}},
\qquad
a_1=\frac{m_2}{m_1+m_2}a.
\]
Combining with Kepler's law
\(4\pi^2a^3=G(m_1+m_2)P^2\) gives the stated function.  Since
\(\sin i\leq1\), the measured \(f\) gives the minimum \(m_2\) for an assumed
\(m_1\).
\end{solution}

\begin{solution}
In the center-of-mass frame,
\[
\bm D_s=\alpha_1m_1\bm x_1+\alpha_2m_2\bm x_2
=\mu(\alpha_1-\alpha_2)\bm r.
\]
Thus equal charge-to-mass ratios eliminate dipole radiation.  For a circular
binary \(P_{\rm dip}\sim v^8\), whereas tensor quadrupole radiation scales as
\(v^{10}\); the dipole is a formally minus-one-PN effect and is strongly
constrained by pulsars.
\end{solution}

\begin{solution}
Assembly from infinity gives
\[
E_g=-\frac{3GM^2}{5R},\qquad
\frac{E_g}{M}=-\frac{3GM}{5R}.
\]
The anomalous acceleration is proportional to
\(\eta_NE_g/M\).  Hence two bodies \(A,B\) in the same external field have
\[
\frac{\Delta a}{g}
=\eta_N\left[\left(\frac{E_g}{M}\right)_A
-\left(\frac{E_g}{M}\right)_B\right],
\]
the Nordtvedt polarization.
\end{solution}

\begin{solution}
For a convention
\[
\dd s^2=-(1+2\Phi)\dd t^2-4\bm A\cdot\dd\bm x\,\dd t
+(1-2\Phi)\dd\bm x^2,
\]
the slow geodesic equation is
\[
\ddot{\bm x}=-\bm\nabla\Phi
+4\bm v\times(\bm\nabla\times\bm A)
\]
up to stationary higher-order terms.  Rescaling the definition of
\(\bm A\) moves the factor \(4\); the metric and observable precession are
convention independent.
\end{solution}
